% ══════════════════════════════════════════
%  CHAPITRE 1 — INTRODUCTION GÉNÉRALE
% ══════════════════════════════════════════
\chapter{Introduction Générale}

\section{Contexte et motivation}
La météorologie joue un rôle crucial dans plusieurs secteurs comme l'agriculture, le tourisme ou la gestion des risques. Actuellement, l'accès à des données météorologiques historiques et en temps réel de manière centralisée reste complexe en Tunisie.

Ce projet vise à résoudre ce problème en utilisant les technologies du Big Data. L'objectif est de créer une plateforme centralisée capable de collecter, traiter et analyser les données météorologiques de plusieurs villes tunisiennes.

\section{Problématique}
Le défi principal de ce projet est de répondre à la question suivante :
Comment concevoir et déployer une plateforme complète capable de :
\begin{itemize}
  \item Collecter des données météo fiables pour l'ensemble des villes tunisiennes ?
  \item Traiter ces flux de données en temps réel ?
  \item Analyser ces données pour en extraire des statistiques utiles ?
  \item Présenter les résultats sur une interface web simple et interactive ?
\end{itemize}

\section{Objectifs du projet}
Pour répondre à cette problématique, nous avons défini les objectifs suivants :
\begin{enumerate}
  \item \textbf{Collecte de données :} Intégrer des APIs météorologiques pour récupérer l'historique et les données actuelles de 221 villes.
  \item \textbf{Infrastructure Big Data :} Mettre en place un système de streaming (Kafka) et une base de données robuste (Supabase).
  \item \textbf{Orchestration :} Automatiser les analyses et la détection d'alertes avec Apache Airflow.
  \item \textbf{Visualisation :} Développer un tableau de bord (Dashboard) interactif pour afficher les données.
\end{enumerate}

\section{Structure du rapport}
Ce rapport s'organise de la manière suivante :
\begin{itemize}
  \item \textbf{Chapitre 2} : Présente l'état de l'art et les outils utilisés.
  \item \textbf{Chapitre 3} : Détaille la méthodologie de travail Scrum.
  \item \textbf{Chapitre 4} : Décrit le Sprint 1 (Infrastructure et Flux temps réel).
  \item \textbf{Chapitre 5} : Décrit le Sprint 2 (Analyses statistiques et Déploiement).
  \item \textbf{Chapitre 6} : Présente les tests globaux et la validation du système.
\end{itemize}
